\chapter{Introduction}\label{ch:1}

Mobility is the property that distinguishes a cellular network from every other access technology, and the handover is the mechanism that delivers it. Each time a user equipment moves beyond the useful coverage of its serving cell, the network must transfer the connection to a neighbour without interrupting the service carried on it. On a dense urban network this happens constantly: across the campaigns reported here a handover command was issued on average once every 11 seconds of driving. Every one of those events is a brief window in which throughput drops and latency rises.

This chapter states the problem the thesis addresses, the reason it is worth addressing, and the objectives against which the work should be judged. Section 1.1 describes how handover is triggered in Long Term Evolution (LTE) and why the triggering rule is structurally incapable of providing advance warning. Section 1.2 quantifies the cost of that limitation on measured data. Section 1.3 states the problem formally, Section 1.4 lists the objectives, Section 1.5 bounds the scope, Section 1.6 summarises the contributions, and Section 1.7 describes the organisation of the remainder of the report.


\section{Background}
\label{sec:1_1}

In an LTE network the user equipment (UE) continuously measures the Reference Signal Received Power (RSRP) and Reference Signal Received Quality (RSRQ) of its serving cell and of the neighbour cells it can detect, and reports those measurements to the network when a condition configured by the network is satisfied \cite{ref1}. The dominant condition governing intra-frequency mobility is Event A3, which is defined in 3GPP TS 36.331 as the condition that a neighbour cell becomes better than the serving cell by a configured offset, and that this relation persists for a configured time-to-trigger (TTT) before the report is sent \cite{ref1}.

Two properties of that definition matter for this thesis. First, the condition is written on a comparison that must already hold: the neighbour must have become better before anything happens. Second, the TTT imposes an additional delay whose purpose is to suppress spurious triggering caused by fast fading. Together they guarantee that by the time a measurement report leaves the handset, the radio environment has already changed. The network then evaluates the report and, if it decides to proceed, issues an RRCConnectionReconfiguration message carrying mobilityControlInfo, which is the handover command itself.


\begin{figure}[htbp]
\centering
\includegraphics[width=0.85\linewidth]{figures/fig01_a3_event.png}
\caption{Event A3 fires only after the neighbour has been better than the serving cell by a configured offset for a configured time-to-trigger. The entry condition is strictly retrospective.}
\label{fig:1_1}
\end{figure}


\section{Motivation}
\label{sec:1_2}

The cost of reacting rather than anticipating is measurable, and it was measured on the data collected for this thesis rather than assumed from the literature. Across four campaigns totalling 2.9 hours of logging, 957 signalling-confirmed handover commands were observed (yielding 938 quality-controlled commands inside retained session blocks, and 934--937 horizon-evaluable events under lead-window boundaries; Sections 3.4 and 3.6), representing a mean interval of 11 seconds with a median gap of 3.7 seconds between consecutive commands. Of those handovers, 26.3\% returned to the cell just vacated within fifteen seconds; each such pair is two signalling exchanges and two service interruptions that produce no net change of serving cell. The signalling log also records 341 RRC connection re-establishment procedures, of which 321 carry the cause reconfigurationFailure and only 3 carry handoverFailure; Section 3.6 treats them accordingly, and they are not described here as radio link failures. Measurement reporting is likewise abundant: 8,136 Event A3 reports were transmitted, which resolve into 2,361 distinct trigger episodes, of which 635 own a handover command once each command is matched to a single episode (Section 5.9).

These are not pathological numbers for a dense urban deployment; they are what the deployed configuration produces. They do, however, establish that even a short prediction horizon has something to act on. A warning far enough ahead of a handover command would leave time for target-cell preparation, and a warning that a handover is likely to be reversed within fifteen seconds would be enough to consider not performing it. Both are stated as conditionals because this thesis does not establish either: Section 5.5 and Appendix A measure what a warning built on this instrument actually delivers at the one-second horizon --- 34\% of commands detected, 179 false-alarm episodes per hour and a median lead of 0.47 s --- and an alarm every 20 seconds with under half a second of warning is not an operating point a radio resource manager can use. The motivation for the task is that the reactive rule leaves a measurable cost on the table; the finding is that one-hertz handset telemetry does not recover enough of it.

The natural question is whether such a predictor already exists. A screening of 108 publications, of which 22 were audited in detail in Chapter 2 and listed in Appendix D, indicates that it does not exist in a form an operator could act upon. Of those 22 models, none holds out the mobility unit --- a whole drive or a whole device session --- when evaluating a per-timestep classification task on measured radio data; none reports a calibration curve, an expected calibration error or a Brier score; and none reports how early its warnings arrive or how many false alarms they generate per unit time. Seven report headline accuracy on tasks whose positive-class prevalence lies between 0.86 \% and 11 \%, where a constant negative prediction already scores between 89 \% and 99 \%. None of them states how the timestamp of a measurement row relates to the interval it summarises, which Chapter 5 shows to be worth more than any modelling choice made in this thesis.


\section{Problem statement}
\label{sec:1_3}

The problem addressed in this thesis is the following. Given the measurements available to a handset at time t --- serving and neighbour signal levels and qualities, their short-term statistics, mobility state and the recent signalling history --- estimate the probability that the network will issue a handover command within h seconds, for several values of h simultaneously, such that:

the estimates are mutually consistent, in the sense that the probability for a longer horizon is never smaller than that for a shorter one;

the estimates are calibrated, in the sense that an asserted probability of p corresponds to an empirical event frequency near p;

the evaluation protocol cannot transfer information from a test drive into training; and

the warning derived from the estimates carries a stated bound on the fraction it will miss on future units, and the unit of that fraction --- a row of the measurement grid, or a handover command --- is stated rather than left to the reader; and

the information available to the model at time t is demonstrably available at time t --- that is, the alignment between the measurement grid and the event clock is verified rather than assumed.

None of these five requirements is satisfied jointly anywhere in the audited literature, and the fourth and fifth are not satisfied anywhere at all in the handover-prediction setting.


\section{Objectives}
\label{sec:1_4}

The specific objectives of this thesis are as follows.

O1. To predict the next handover at five look-ahead times --- 0.5, 1, 2, 3 and 5 seconds --- using only quantities observable at the handset at the instant of prediction.

O2. To establish an evaluation protocol that cannot leak information between training and test, by grouping every split at the level of the whole drive, and to measure how much that protocol is worth relative to the random-row splitting that is standard in the comparator literature.

O3. To make the resulting probabilities usable by an operator: mutually coherent across horizons, calibrated in value, and accompanied by a distribution-free bound on the per-drive miss rate.

O4. To test whether the result generalises beyond the campaigns on which it was developed, both to a corridor and speed regime held out for model fitting (acknowledging that vehicle speed and geographic corridor are structurally collinear and confounded in observational drive data, and that Campaign 4 was evaluated in retrospective diagnostic exploration of the lag repair) and to an external public dataset collected on the same operator by an earlier cohort in our department [38] (evaluating cross-dataset generalisation on a common signalling schema).

O5. To explain the mechanism underlying the prediction, rather than reporting accuracy alone, and in doing so to characterise the behaviour of the deployed mobility configuration on this network.


\section{Scope and limitations}
\label{sec:1_5}

The scope of this work is deliberately bounded, and the bounds are stated here rather than deferred to the conclusion.

The study covers one commercial LTE operator, four measurement campaigns, two mobility regimes and 57 drives. It does not claim that the numerical results transfer to another operator or another radio access technology, although Chapter 5 reports an external validation that bears on that question. The measurement export is sampled at one hertz. At the raw inter-event interval level, 90.5\% of consecutive handover events have a separation $\ge 1\text{ s}$. However, when an operational 2~s post-handover blanking window is enforced to suppress immediate ping-pong contamination, the effective fraction of handover commands with an eligible prediction row in their 1~s lead window is capped at 64.2\% (with 35.8\% of events falling within the blanking window of a preceding command); this is an intrinsic property of the physical logging rate coupled with the operational blanking policy, not of the predictive models.

The system developed here is an offline predictor evaluated on recorded data. It is not integrated into a network, and no claim is made that deploying it would improve any network key performance indicator. In particular, no causal claim is made about the benefit of acting on the warning, for the reason given in Chapter 5: the logging policy is deterministic, so the propensities required for off-policy evaluation are zero or one and the causal estimand is not identified on observational data of this kind. Only a counting upper bound is reported.

Finally, the thesis makes no claim relating to human health. No human subjects were involved, no exposure measurements were taken, and the impact analysis in Chapter 6 is confined to what the measurements support.


\section{Contributions}
\label{sec:1_6}

The contributions of this thesis are the following six, each supported by a specific result in Chapter 5. They are stated as they stand after the alignment audit of Section 5.1, which revised two of them downwards from an earlier draft of this work.

1. A timestamp-alignment audit for drive-test prediction studies, a measurement of what it is worth, and an identification of the columns that carry it. In 76\% of handovers the one-hertz export row stamped before the handover command already carries the post-handover serving cell; removing that information by lagging every export feature one row moves the one-second AUPRC from 0.600 to 0.179 and single-feature dwell AUROC from 0.873 to 0.664. While Campaign 4 serves as an out-of-sample holdout for model weights and hyperparameters, the alignment defect and one-row lag repair were derived through retrospective exploratory analysis across the initial campaigns; hence Campaign 4 is not a blind holdout for pipeline feature engineering. Lagging export features simultaneously eliminates post-handover cell identity contamination and introduces one second of information ageing; while lagging the non-corrupted GPS mobility block incurs negligible additional cost (Table~\ref{tab:5_2b}), within the serving-radio features both effects are confounded at 1~Hz, rendering the one-row lag a conservative safeguard against look-ahead bias under assumed export semantics. The reduction is concentrated in the serving-cell radio scalars (accounting for 96\% of the fall against 2\% for the neighbour and gap columns), because columns such as serving\_rsrp report whichever cell the export treats as serving and change their referent at the flip. The audit is three lines of code and is reproducible on any export that carries a serving-cell column.

2. An account of why that happens, tested against the predictions it makes. A row written at the end of its second carries the post-handover cell exactly when the handover completes before the second closes. The prediction is supported across instruments: the defect appears on 4,712 handovers from 3 operators in an external multi-carrier XCAL corpus [40] when that corpus is downsampled to the end of each second (93.2\%) and collapses on the same handovers when downsampled to the start (4.1\%), while being absent from an instantaneous sampling tool [37]. On our own captures, an in-sample diagnostic check on the 15 handovers whose completion crosses the second boundary confirms zero contamination on the row stamped before the command (0/15; 95\% Wilson score interval [0.0\%, 20.4\%], with the remaining 211 non-straddling residual contaminated cases explained by a 1-row exporter refresh buffer delay). The alignment error is therefore a property of the aggregation rule, not of a vendor or a network --- and, read the other way, the contamination curve near the end of the second estimates an inferred logging and periodic export commit characteristic at roughly 60 ms without any sub-second logging.

3. A measurement of what the unit of splitting is worth on drive-test radio data, on a ladder of four protocols rather than the usual two. Against leave-one-campaign-out, splitting by random 180-second chunks of the same continuous session inflates the one-second AUPRC by 13 \% and splitting rows at random by 48 \%; the inflation is architecture-dependent, so the protocol determines the ranking of models and not only their level.

4. A discrete-time hazard formulation of handover prediction that renders five look-ahead horizons mutually consistent from a single model fit, for any learner, without a held-out calibration split and without post-hoc correction. Independent per-horizon classifiers violate ordering on 40.2\% of rows; the hazard model cannot. The claim is deliberately narrow in two directions. Time-to-event modelling of cellular mobility is not new --- Section 2.5 places this work against the cell-sojourn-time and handover-rate literature, which derives analytical marginal distributions under assumed mobility and geometry --- and what is offered here is a learned, covariate-conditional, per-sample hazard fitted on measured drive-test data and scored against decoded ground truth. And the comparison is reported against the repairs a reviewer would ask for --- a cumulative-maximum projection, and isotonic calibration composed with it --- which also reach zero violations. After those concessions the argument for the formulation is that it reaches coherence and competitive calibration from one fit with no held-out calibration split; it does not rank better than the independent control, and the censoring-aware likelihood it inherits is inert on data in which no row is censored inside the horizon grid.

5. A boundary result on distribution-free risk control for handover warning, reported as such rather than as a delivered guarantee. Both the exchangeable unit and the estimand are stated and tested. The unit is tested in the regime where exchangeability holds and the regime where it does not: calibration and test units drawn from the same pool of chunks, and calibration on one campaign with the guarantee claimed for another. The estimand is carried through in two forms --- a bound over rows and a bound over handover commands --- because they differ, and the difference produces the more useful result: at the one-second horizon 35.8\% of commands have no usable row in their lead window under the post-handover blanking policy and are missed at every threshold, so no per-handover bound is expressible on a one-hertz grid at that horizon without continuous alarm. Two floors therefore bind: the number of calibration units, which is a direct corollary of the conformal risk control bound \cite{ref21}, and the preprocessing blanking window. The summary of this contribution highlights these two limits: the empirical mean miss loss across the 12 pairwise campaign rotations was 0.153 at a target of 0.20, but cross-campaign coverage cannot be formally certified because campaign exchangeability is violated by macro-environmental distribution shift (and pairwise rotation over 4 campaigns yields effective sample size $n=4$ pseudoreplication). And the estimand an operator cares about is not expressible at all at the shortest horizon. What this thesis contributes is therefore the two limits and the measurements that establish them, not a certified warning system.

6. Two corrected characterisations of the deployed network that the prediction literature has not modelled. The first is that the A3 report-conversion rate has no unit-free value: it reads 63\% declined per transmitted report, 30\% per trigger episode when overlapping episodes may claim the same command, and 73\% once each command is matched to exactly one episode. The loose episode rule records 1,647 episode-to-command claims over a corpus containing only 957 commands, an over-counting factor of 2.6. What survives every unit, and sharpens under the strictest, is the separation between the mobility configuration and the monitoring configurations (37\% declined against 95\%). The second is the deployed offset distribution weighted by handovers rather than by configurations (75\% of A3-triggered handovers fire under a positive offset).


\section{Organisation of the report}
\label{sec:1_7}

The remainder of this report is organised as follows. Chapter 2 reviews the four literatures that bear on this problem --- cellular measurement studies, handover prediction, survival analysis and distribution-free uncertainty quantification --- presents a protocol audit of the 22 most comparable prediction models, and sets out the methodology adopted, including the approaches that were considered and rejected. Chapter 3 describes the measurement campaign, the signalling-decoding procedure and the construction of the dataset. Chapter 4 develops the predictive formulation, the feature construction, the models compared, the evaluation protocol and the risk-control procedure. Chapter 5 reports and discusses the results against the objectives of Section 1.4. Chapter 6 demonstrates how the work addresses the course outcomes, program outcomes, knowledge profiles and complex engineering problem and activity attributes required of a capstone project (omitted when compiled in journal mode). Chapter 7 concludes and identifies the work that follows from the limits of the present study.

